%% notes-tex/modeling/scaling-laws/sections/5-torchcell.tex

\section{What an ablation looks like on the trigenic build\secstatus{todo}}
\label{sec:torchcell}

\begin{keybox}
\textbf{Measured.} No number here is a torchcell result. The constants
are the published ones and panel~i is synthetic.

\textbf{Not measured.} Every number a scaling section of the manuscript or thesis
would state: the exponents, the floor, and whether a curve fit on small runs
predicts a large one. What follows is the design, with the choices each sweep
forces written down so a run can be launched against it.
\end{keybox}

\notesec{The axes available}
The trigenic build of experiment 025 holds 376{,}732 records with bit-identical
labels to the 010 build, organized as 420 recurring query pairs plus one array
gene each, and it already carries the two splits an ablation needs: the published
random-over-records split and a query-pair-disjoint split under which no query
double is shared between parts. The equivariant cell graph transformer trained on
it is one architecture family whose width and depth can be scaled together. The
axes are therefore available now; what is missing is the sweep.

\notesec{The data sweep first}
Data is the binding axis here, so the $D$-sweep is the one to run before any
other. Nested subsets of the disjoint-split training set at 1/64, 1/16, 1/4 and 1,
drawn by query pair so that a smaller subset removes whole screens, against the
one fixed disjoint test set; a model at the current configuration's size; the
loss recorded per epoch and the validation-selected checkpoint scored on test.
Four points give a slope and, with the held-out largest subset, the first
extrapolation test. The reading the sweep is for: whether the loss is still
falling with $D$ at the full set, which says a new screen is worth more than a
larger model, or has flattened, which says the opposite.

\notesec{The parameter sweep}
Five sizes of the transformer, width and depth in the fixed ratio of the current
configuration, spaced evenly in $\log N$ over about two decades, with the gene
embedding table counted separately and the curve fit on the rest; three seeds at
the two smallest sizes; learning rate tuned per size, or transferred under $\mu$P
once the parametrization is checked against the model code. Every run on the full
disjoint training set, so that this sweep reads $L(N)$ with data as fixed as it
can be made.

\notesec{The floor from the replicates}
Before either sweep, the replicate standard deviations of the test records give
the noise-variance bound on $E$ in the loss's units. Whether the trigenic
interaction score's released uncertainty is a sample standard deviation over
colonies or a bootstrap standard error, and over how many colonies, is already
recorded for the Kuzmin and Costanzo screens in the noise-computation note, and
the bound follows from it. That number is the first thing to set beside a fitted
$E$.

\notesec{Compute}
The cost per record of the transformer depends on how many gene tokens a record
carries and on the graph penalty, so $C$ is read from a profiler at each size
rather than from $6ND$. The IsoFLOP sweep is third in priority: it answers how to
split a budget, which matters once the first two sweeps have shown that the
budget buys anything.

\notesec{What the section in the manuscript would then say}
A scaling section is one figure, three panels, $L$ against $D$, against $N$ and
against $C$, each with its fitted exponent and interval, the held-out largest run
marked, and $E$ drawn against the replicate floor. Until the sweeps have run, that
figure is a plan, and the sweeps above are its specification.
